\section*{Chapter \ref{ch:m2:swap-clearing-and-the-ccp-basis} --- Swap Clearing and the Clearing-House Basis}

\begin{solution}{exo:m2:swap-clearing-and-the-ccp-basis:1}
$36.4 \times 81\,109 = \text{USD}~2.95$ million: 36.4 basis points of rate, applied to
the DV01.
\end{solution}

\begin{solution}{exo:m2:swap-clearing-and-the-ccp-basis:2}
A receiver loses when rates rise: $3 \times 81\,109 = \text{USD}~243\,327$, paid in cash
through its clearing broker to the house the same day or the next morning.
\end{solution}

\begin{solution}{exo:m2:swap-clearing-and-the-ccp-basis:3}
If \omterm{def:m2:swap-clearing-and-the-ccp-basis:mandate}{porting} succeeds, the client's swaps and its margin move to another
clearing member and continue unchanged: its hedges survive. If it fails, the
house closes the positions out, returning the client's margin net of losses;
the client must re-enter its hedges in whatever market prevails, having been
unhedged in between.
\end{solution}

\begin{solution}{exo:m2:swap-clearing-and-the-ccp-basis:4}
USD~76\,442 over the swap's life; divided by the DV01 of 81\,109, a running basis
of 0.94 basis points a year.
\end{solution}

\begin{solution}{exo:m2:swap-clearing-and-the-ccp-basis:5}
At the house where clients pay fixed. Dealers who receive there end up one way
and must hedge elsewhere, posting margin twice; they ask a higher fixed rate to
receive, so the par rate there is higher. The analysis of 2017 found swaps at
CME quoted above the same swaps at LCH for exactly this reason.
\end{solution}

\begin{solution}{exo:m2:swap-clearing-and-the-ccp-basis:6}
Yes: EUR~9 billion is above the EUR~8 billion threshold of the final phase.
Initial margin is exchanged only above the EUR~50 million threshold per
counterparty group: EUR~20 million, if nothing else in the group uses the
threshold.
\end{solution}

\begin{solution}{exo:m2:swap-clearing-and-the-ccp-basis:7}
1.15, 2.30 and 3.45 basis points per one-way position. At 50, two houses need
4.60, against the 3.4 of the 2017 example: the same order of magnitude, the
difference coming from the illustrative volatility, margin period and funding
spread, and from real margin models that net across a dealer's portfolio.
\end{solution}

\begin{solution}{exo:m2:swap-clearing-and-the-ccp-basis:8}
The two swaps offset in risk but not in margin: each house margins its side on
its own, so the ``arbitrage'' posts two initial margins for thirty years. At the
market's funding costs that is what the 3 basis points pay for. Only a firm
whose existing positions make each new trade \emph{reduce} its margin at both
houses earns the spread, and that is the business dealers compete for.
\end{solution}

\begin{solution}{pb:m2:swap-clearing-and-the-ccp-basis:1}
\textbf{1.} $2 \times 10^9 \times 8.111 \times 10^{-4} = \text{USD}~1.62$ million per basis point.
\textbf{2.} $36.4 \times 1.622 \approx \text{USD}~59.1$ million.
\textbf{3.} $5 \times 1.622 = \text{USD}~8.11$ million, paid by the dealer, who receives
fixed at A; at B, where it pays fixed, it is paid the same amount.
\textbf{4.} Each house guarantees only its own trades and must be able to close
out a defaulter's portfolio \emph{there}; a hedge at another house does not
reduce its loss.
\textbf{5.} It would fall towards zero: a flat position at A has no market risk
there, and only a small margin for basis and liquidity.
\textbf{6.} The same USD~59.1 million.
\textbf{7.} USD~1.53 million each, USD~3.06 million for both.
\textbf{8.} 0.94 basis points for one, 1.88 for both.
\textbf{9.} 4.60 basis points for both.
\textbf{10.} $\sqrt 2$ times as much: USD~83.5 million at each house.
\textbf{11.} By paying fixed at A to a dealer who receives there, its position at
A is flat: no significant initial margin at A, and none at B, where it no longer
needs a hedge.
\textbf{12.} Up to 1.88 basis points of rate: what the two one-way margins would
cost over the ten years.
\textbf{13.} Because every dealer at A faces the same clients: all are receivers
there, and none wants to add to the imbalance by receiving more.
\textbf{14.} In the fixed rate they pay at A, higher than at B by about the cost
of the dealers' two margins.
\textbf{15.} Clients who receive fixed at A (pension funds, insurers), a rule
moving flows between houses, or a change in either house's margin model or
cross-margining with the other.
\textbf{16.} Because it must be funded: the house pays on cash collateral less
than the firm pays to borrow it, and securities posted must be financed in
repo; the gap is the funding spread.
\textbf{17.} They use historical simulation of whole portfolios, add-ons for
liquidity and concentration, and floors; netting within a portfolio makes the
incremental margin of a trade depend on everything else the member holds.
\textbf{18.} It forces flows to a house regardless of its portfolio, creating
one-way books where there were none, and so creating or widening a basis
between that house and the others.
\textbf{19.} Named result: \emph{the break-even basis} for ten-year swaps is 1.88 basis
points: the running cost of margining one-way positions at both houses, which a
dealer should be willing to pay to net at house A instead.
\textbf{20.} Because margin is computed at each clearing house separately, and
what it costs to hold a swap depends on who else trades there.
\end{solution}

\begin{solution}{iq:m2:swap-clearing-and-the-ccp-basis:1}
Bilateral swaps had created a web of exposures between large banks, opaque to
regulators and to each other, which in 2008 amplified the fear of any one
default. The G20 required standard swaps to be cleared, reported and, where
possible, traded on platforms. For a dealer: its counterparty on most swaps
is a clearing house; it posts initial and variation margin daily; its
balance sheet benefits from netting; its costs now include margin funding and
default-fund contributions.
\iqlookfor{the counterparty web as the problem and margin as the new cost.}
\end{solution}

\begin{solution}{iq:m2:swap-clearing-and-the-ccp-basis:2}
Variation margin settles the daily change in a position's value, so that
neither side accumulates an unpaid gain or loss. Initial margin is collateral
posted up front, and kept, to cover the loss the house or counterparty could
suffer while closing out a defaulter's positions over the margin period of
risk; it is returned when the position is closed.
\iqlookfor{settlement of past moves versus buffer for future ones.}
\end{solution}

\begin{solution}{iq:m2:swap-clearing-and-the-ccp-basis:3}
The difference in par rate between identical swaps cleared at CME and at LCH.
Client flows at one house are one-directional, dealers end up one way at
each house and post two margins; the basis prices that funding cost. It is not
an arbitrage: capturing it recreates the two margins it pays for, unless your
existing positions make the trades margin-reducing.
\iqlookfor{the flow imbalance and margin mechanism, and why it is not free
money.}
\end{solution}

\begin{solution}{iq:m2:swap-clearing-and-the-ccp-basis:4}
Margin valuation adjustment: the present value of funding the initial margin a
trade requires over its life. Estimate it by projecting the trade's
incremental initial margin through time (from its future DV01 profile, or by
simulating the portfolio margin), multiplying by the funding spread and
discounting; express it as a running spread by dividing by the trade's DV01.
\iqlookfor{incremental, not standalone, margin, and a projection through
time.}
\end{solution}

\begin{solution}{iq:m2:swap-clearing-and-the-ccp-basis:5}
Bilateral: the dealer is the counterparty, so credit risk is managed through the
CSA's collateral terms; if both parties are above the thresholds, initial margin
is exchanged and segregated, usually larger than at a clearing house because
the margin period is longer (ten days rather than five) and netting narrower;
the dealer charges for credit and funding (the valuation adjustments of One
Quant Book 6). The client gains flexibility in terms and avoids clearing fees
and default-fund exposure through its broker.
\iqlookfor{credit and margin differences, and their pricing.}
\end{solution}

\begin{solution}{iq:m2:swap-clearing-and-the-ccp-basis:6}
Replicate the house's published margin methodology: historical scenarios
applied to the member's portfolio, with its liquidity and concentration
add-ons. Precompute the portfolio's scenario P\&L vector each morning; for a new
trade compute its scenario P\&L vector (linear in sensitivities for speed, full
revaluation for large trades), add it, and read the new quantile; the difference
is the incremental margin. Keep the vectors in memory, update them on each
trade, and test against the house's own margin calculator.
\iqlookfor{scenario vectors, incremental quantiles, and validation against the
house.}
\end{solution}
